\documentclass[xcolor={dvipsnames}]{beamer}

\usepackage{../assets/pres-template_MOW}

%--------------------------------------------------------------------------
% Specific to this document ---------------------------------------
%--------------------------------------------------------------------------
\setbeamercovered{transparent}

%%%%%%%%%%%%%%%%%%%%%%%%%%%%%%%%%%%
\setlength{\tabcolsep}{1.3pt}
\title{Social Science Inquiry II}
\subtitle{Week 1: Course introduction}
\date{Winter 2024}
\author{Molly Offer-Westort}
\institute{Department of Political Science, \\University of Chicago}


\usepackage{Sweave}
\begin{document}
\input{slides_11-concordance}

%-------------------------------------------------------------------------------%
\frame{\titlepage
\thispagestyle{empty}
}
%-------------------------------------------------------------------------------%
\begin{frame}{Overview}

This course:\pause
\begin{itemize}
\item Part of a sequence \pause
\begin{itemize}
\item Fall: research design \pause
\item Winter: methods \pause
\item Spring: practical applications
\end{itemize}
\end{itemize}

\pause

Did you take SOSC 13100 in fall quarter of 2023?

\end{frame}


%%%%%NOTE%%%%%
\note{
\scriptsize \singlespacing

\begin{itemize}
\item xxxx
\end{itemize}

}

%-------------------------------------------------------------------------------%
\begin{frame}

\begin{center}
Main take-aways from fall quarter?
\end{center}

\end{frame}


%%%%%NOTE%%%%%
\note{
\scriptsize \singlespacing

\begin{itemize}
\item xxxx
\end{itemize}

}

%-------------------------------------------------------------------------------%

\begin{frame}{Course objectives}
\pause
\begin{itemize}
\item Students will approach data thinking like a social scientist; i.e., thinking about what data inputs are, how they’re measured, relationships among variables, thinking about data generating processes when making inferences\pause
\item Build on foundations of descriptive and causal inference introduced in fall SSI I, with respect to both experimental and observational data \pause
\item Have facility with basic statistical tests and methods \pause
\item Have facility with basic statistical software programming and data visualization 
\end{itemize}

\end{frame}


%%%%%NOTE%%%%%
\note{
\scriptsize \singlespacing

\begin{itemize}
\item xxxx
\end{itemize}

}

%-------------------------------------------------------------------------------%
%%% [[TEMPLATE]] %%%
\begin{frame}{Course objectives}

By the end of the quarter, you should be able to \dots\pause
\begin{itemize}
\item Formulate a falsifiable research ``hypothesis'' \pause
\item Determine what evidence (data, statistical tests) you would need to falsify that hypothesis \pause
\item Interpret results from statistical analysis \pause
\item Run basic statistical analysis in R
\end{itemize}

\end{frame}


%%%%%NOTE%%%%%
\note{
\scriptsize \singlespacing

\begin{itemize}
\item xxxx
\end{itemize}

}

%-------------------------------------------------------------------------------%
\begin{frame}{Overview}

What about qualitative research methods? \pause
\begin{itemize}
\item Data can be both quantitative and qualitative. \pause
\item Principles of research design can be shared across methodological approaches. 
\end{itemize}

\end{frame}


%%%%%NOTE%%%%%
\note{
\scriptsize \singlespacing

\begin{itemize}
\item xxxx
\end{itemize}

}

%-------------------------------------------------------------------------------%
\section{Overview}
\begin{frame}{Overview}

\begin{itemize}
\item High level: how to answer questions with data \pause
\item Summarizing univariate and multivariate data numerically and visually, data exploration \pause
\item Probability as a model of the world \pause
\item Formalizing our uncertainty about that model, using statistics \pause
\item What it requires to infer causal relationships \pause
\item Linear regression \pause (and inference with) \pause
\item Beyond linear regression: other statistical methods
\end{itemize}

\end{frame}


%%%%%NOTE%%%%%
\note{
\scriptsize \singlespacing

\begin{itemize}
\item xxxx
\end{itemize}

}

%-------------------------------------------------------------------------------%

\begin{frame}{What we do with data}

We'll get started on working with data, and we'll also get started on thinking critically about how you use data to answer questions.\\~\ \pause


What data would you need to make the argument in the article below?

\begin{figure}[htbp] %  figure placement: here, top, bottom, or page
   \centering
   \includegraphics[width=.8\textwidth]{../assets/guardian.png} 
   % \caption{example caption}
   % \label{fig:example}
\end{figure}

\end{frame}


%%%%%NOTE%%%%%
\note{
\scriptsize \singlespacing

\begin{itemize}
\item xxxx
\end{itemize}

}

%-------------------------------------------------------------------------------%
\begin{frame}{Inferential questions}

\begin{itemize}
\item What can the data you \textit{do} have tell you about data you \textit{don't} have? \pause
\item What data would you need to answer questions about \textit{what would have happened}? \pause
\item What can we say about our \textit{uncertainty} about estimates or predictions?
\end{itemize}

\end{frame}


%%%%%NOTE%%%%%
\note{
\scriptsize \singlespacing

\begin{itemize}
\item xxxx
\end{itemize}

}

%-------------------------------------------------------------------------------%
\begin{frame}{Statistical software}

Programming in R.\pause
\begin{itemize}
\item Assuming you're starting from scratch. \pause
\item But if you do have some programming skills, with statistical software, python, or other languages, there can be a lot of carryover. \pause
\item If you don't have any baseline skills, focus on the rstudio primers \url{https://rstudio.cloud/learn/primers}
\end{itemize}

\end{frame}


%%%%%NOTE%%%%%
\note{
\scriptsize \singlespacing

\begin{itemize}
\item xxxx
\end{itemize}

}

%-------------------------------------------------------------------------------%
\begin{frame}{Reference texts}

\begin{center}
{[}See \href{https://docs.google.com/document/d/1wiQwd5HgVohVTC46dfHo67dQ2KBJjNmdO2Pyv5e2Q_s/edit?usp=sharing}{syllabus} {]}
\end{center}

\end{frame}


%%%%%NOTE%%%%%
\note{
\scriptsize \singlespacing

\begin{itemize}
\item xxxx
\end{itemize}

}

%-------------------------------------------------------------------------------%

\begin{frame}{Asking and answering questions}

\begin{itemize}
\item Google! \pause
\item Ask on our course StackOverflow team: \url{https://stackoverflowteams.com/c/sosc13200-w24-6} \pause {[}I'll send invites shortly{]}  \pause
\item Use general StackOverflow (\url{https://stackoverflow.com/}) for \texttt{R} questions or CrossValidated (\url{https://stats.stackexchange.com/}) for stats. \pause
\end{itemize}

Class participation includes asking and \textit{answering} questions on the class StackOverflow.

\end{frame}


%%%%%NOTE%%%%%
\note{
\scriptsize \singlespacing

\begin{itemize}
\item xxxx
\end{itemize}

}

% %-------------------------------------------------------------------------------%

\begin{frame}{Assessment}

\begin{itemize}
\item Homework: 60\%
\item Final project: 30\%
\item Participation: 10\%
\end{itemize}

\end{frame}


%%%%%NOTE%%%%%
\note{
\scriptsize \singlespacing

\begin{itemize}
\item xxxx
\end{itemize}

}

%-------------------------------------------------------------------------------%

\begin{frame}{Homework}

\begin{itemize}
\item Solution sets will be posted \textit{at the same time} as problem sets. \pause
\item Do as much as you can on the problem set before checking the solutions. \pause
\item Check your work, and then fill out a form on how you did, what you understood and didn't.\pause
\item You get marked both on completion of the problem set, \textbf{AND} filling out the form.\pause
\item (If you find errors in the solution set, post them on the class StackOverflow and you will get extra credit) \pause
\item For homework assignments, always submit \textit{both} your .R file showing your work, and and a compiled .pdf\ file on Canvas.
\end{itemize}

\end{frame}


%%%%%NOTE%%%%%
\note{
\scriptsize \singlespacing

\begin{itemize}
\item xxxx
\end{itemize}

}

%-------------------------------------------------------------------------------%

\begin{frame}{Homework grading}

check(+/-) \pause
\begin{itemize}
\item Check: You fully completed the assignment, and submitted all components. (A) \pause
\item Check plus: You went above and beyond, your solutions were clear and detailed. (A+)  \pause
\item Check minus: You made an attempt, but it wasn't complete. Maybe you didn't submit all components, or didn't fully answer some of the questions. (B or C) \pause
\item Unmarked: You did not submit enough of an assignment for credit. 
\end{itemize}

\end{frame}


%%%%%NOTE%%%%%
\note{
\scriptsize \singlespacing

\begin{itemize}
\item xxxx
\end{itemize}

}

%-------------------------------------------------------------------------------%


\begin{frame}{Assignment for Monday}

\begin{itemize}
\item Download software on your own computer. \pause
\item Set up your working directory. \pause
\item Compile and submit a report.\pause
\item (Fill out the feedback form as well, even though there aren't solutions.) 
\end{itemize}

\end{frame}


%%%%%NOTE%%%%%
\note{
\scriptsize \singlespacing

\begin{itemize}
\item xxxx
\end{itemize}

}

%-------------------------------------------------------------------------------%
\begin{frame}{Final Project}

\pause
\begin{itemize}
\item More in-depth, open-ended problem set. \pause
\item Make sure you understand the homework solutions, the final will not have a solution set. \pause
\item You pick from approved data sets--or propose your own. \pause
\item Run through a series of coding exercises and statistical analyses. 
\end{itemize}

\end{frame}


%%%%%NOTE%%%%%
\note{
\scriptsize \singlespacing

\begin{itemize}
\item xxxx
\end{itemize}

}

%-------------------------------------------------------------------------------%
\begin{frame}{Github}

\begin{itemize}
\item \url{https://github.com/UChicago-pol-methods/SOSC13200-W24}\pause
\item Serves as a repository for class documents. \pause
\item You don't need a GitHub account, or need to know how to use git. \pause
\item But if you \textit{do} find mistakes in slides or homeworks, you can get extra extra credit for submitting a pull request on GitHub with revisions. 
\end{itemize}

\end{frame}


%%%%%NOTE%%%%%
\note{
\scriptsize \singlespacing

\begin{itemize}
\item xxxx
\end{itemize}

}

%-------------------------------------------------------------------------------%
\begin{frame}

\begin{center}
Other questions/concerns about class policies?
\end{center}

\end{frame}


%%%%%NOTE%%%%%
\note{
\scriptsize \singlespacing

\begin{itemize}
\item xxxx
\end{itemize}

}

%-------------------------------------------------------------------------------%
\begin{frame}{Getting started in R}

\begin{center}
{[} R script {]}
\end{center}

\end{frame}


%%%%%NOTE%%%%%
\note{
\scriptsize \singlespacing

\begin{itemize}
\item xxxx
\end{itemize}

}


%-------------------------------------------------------------------------------%
\end{document}
%
%-------------------------------------------------------------------------------%

%%% [[TEMPLATE]] %%%
\begin{frame}{Frametitle}

\begin{itemize}
\item xxx
\end{itemize}

\end{frame}


%%%%%NOTE%%%%%
\note{
\scriptsize \singlespacing

\begin{itemize}
\item xxxx
\end{itemize}

}

%-------------------------------------------------------------------------------%
